\section{Introduction}
\label{sec:introduction}

\subsection{Motivation and context}

Autonomous vehicles approaching an unsignalised give-way junction must reason
about another road user's future motion while trading progress against risk.
In a modular stack, the realised manoeuvre is not determined by the predictor
alone: its distribution is consumed by a risk-allocation rule and trajectory
solver, may be overridden by a safety supervisor, and changes the next
observation. Consequently, a model can improve displacement error or
likelihood without producing a uniform closed-loop benefit.

This dissertation began with the planning-level hypothesis that adaptive risk
allocation would outperform fixed risk in the give-way scenario. Pilot
analysis showed that a shared supervisor and predictor could compress the
observable difference between planning policies. Rather than treating this as
a failed comparison, the final study follows the mechanism upstream: it first
tests whether scenario-specific prediction can be improved, then whether
explicit temporal interaction modelling adds value, and finally whether the
frozen predictor improvement changes the closed-loop safety--efficiency
frontier.

\subsection{Research gap}

The dissertation addresses three linked gaps. First, motion-prediction studies
often stop at open-loop ADE, FDE or likelihood, leaving the downstream value of
an improvement uncertain. Second, architectural complexity and attention are
frequently compared without a simple task-adaptation baseline and an
approximately capacity-matched non-attention control. Third, a claim that an
adaptive risk rule is superior can be an artefact of comparing it with a
single fixed setting rather than multiple prespecified fixed operating points
and of omitting moderators such as predictor and target behaviour.

\TODO{Support each gap with a critical synthesis of primary literature. State
what each cited paper evaluates, what it leaves unresolved, and why the three
gaps intersect in this system.}

\paragraph{Central thesis.}
In this controlled give-way setting, task adaptation rather than additional
interaction-model complexity provides the reliable offline prediction gain;
after that gain passes through calibration, risk allocation, optimisation and
supervision, its closed-loop value becomes context dependent. Consequently,
neither better offline prediction nor adaptive risk alone establishes uniform
system-level superiority.

\subsection{Research questions and hypotheses}

The study asks:

\begin{description}
  \item[RQ1] Does task adaptation improve in-distribution give-way trajectory
  prediction relative to a pretrained MultiPath control?
  \item[RQ2] Do explicit interaction sequences and Transformer encoders add
  benefit beyond simple adaptation and approximately matched MLP capacity?
  \item[RQ3] Do offline prediction gains translate consistently into
  closed-loop benefit across risk policies and target styles?
  \item[RQ4] Does adaptive risk allocation dominate the fixed-risk
  safety--efficiency frontier?
\end{description}

Each research question maps to one core hypothesis: task adaptation improves
prediction relative to the pretrained control (H1); explicit-sequence
Transformer adaptation adds benefit beyond simple task adaptation that is not
explained by matched non-attention capacity (H2);
the offline gain transfers consistently to closed-loop safety--efficiency
outcomes (H3); and adaptive risk dominates the fixed-risk frontier (H4).
Sequence ablation, calibration-tail analysis, deployment gates and
collision-filtered retraining are mechanism or robustness checks, not
additional headline hypotheses. Table~\ref{tab:hypotheses} reports the four
evidence-bounded verdicts.

\subsection{Contributions}

The dissertation makes four focused contributions:

\begin{enumerate}
  \item evidence that simple task adaptation, rather than the tested increase
  in interaction-model complexity, is the reliable source of offline gain;
  \item mechanism diagnostics showing that sequence use and aggregate
  calibration quality are insufficient to establish model superiority;
  \item a frozen closed-loop factorial evaluation demonstrating that predictor
  value is conditional on risk policy and target style; and
  \item a frontier-based re-evaluation of the original adaptive-risk claim,
  supported by rollout-disjoint data, frozen selection and auditable
  robustness checks.
\end{enumerate}

The contribution is not a claim of general Transformer superiority or
real-road safety. It is an empirical account of where task adaptation helps,
where added sequence complexity does not, and why predictor quality and risk
allocation must be assessed as a coupled system.

\subsection{Paper organisation}

Section~\ref{sec:related} critically reviews prediction, interaction modelling
and risk-aware planning. Section~\ref{sec:problem} defines the coupled problem.
Sections~\ref{sec:methods} and~\ref{sec:experiments} specify the reproducible
protocol, Section~\ref{sec:results} reports the evidence, and
Sections~\ref{sec:discussion}--\ref{sec:conclusion} interpret the findings and
their boundaries.
